\chapter{Il lemma di Yoneda e le categorie di prefasci}\label{cap_yoneda}
\subsubsection*{Esercizi}
\begin{enumerate}
 \item
 \item
 \item
 \item
 \item
\end{enumerate}


\section{Il lemma di Yoneda}

blabla

\section{I prefasci come colimiti di rappresentabili}\label{sec_colim_rappr}

Il risultato fondamentale di questa sezione è \ref{prop_colim_rappr}, che costituisce il punto di partenza per dimostrare una proprietà universale delle categorie di prefasci. Ne dedurremo anche un teorema di caratterizzazione (si veda \ref{}) e una proprietà di stabilità (si veda \ref{}) per le categorie di prefasci. Si ricordi (il duale del)la definizione \ref{elts_F} di categoria degli elementi di un funtore \(A\colon \ctD \fun \ctSet\): $\Elts\ctD A$ ha
\begin{itemize}
      \item per oggetti le coppie \((X \in \ctD, x \in AX)\);
      \item per frecce \(f \colon (X,x) \to (Y,y)\) le frecce \(f \colon Y \to X\) in \(\ctA\) tali che \(A(f)(y)=x\).
\end{itemize}
Si osservi che la categoria $\Elts\ctD A$ è dotata di un ovvio funtore dimenticante controvariante
 \[\dmFun{\Phi_{A}}{\Elts\ctD A}{\ctD^\op}\]
 che manda $f \colon (X,x) \to (Y,y)$ in $f \colon Y \to X$.

Ciò che ora dimostriamo che è usando questa `costruzione degli elementi', ogni funtore \(A \colon \ctD \fun \ctSet\) si può esprimere in modo canonico come colimite
di funtori rappresentabili.

\begin{proposition}\label{prop_colim_rappr}
 Sia \(\ctD\) una categoria piccola; si consideri la composizione di funtori 
 \[\xymatrix{
    \Elts\ctD A \ar[r]^-{\Phi_A} & \ctD^\op \ar[r]^-{\coyon_\ctD} & \copresh\ctD.
 }\]
 Allora, per ogni funtore \(A \colon \ctD \fun \ctSet\) si ha
 \[\colim \big(\coyon_{\ctD} \cdot \Phi_{A}\big) = \left\langle A, \{\xi_x \colon \ctD(X,\blank) \nat A\}_{(X,x) \in \Elts\ctD A}\right\rangle\]
 dove \(\xi_x\) è la trasformazione naturale che corrisponde all'elemento \(x \in AX\) nella biiezione naturale
 \(\Nat(\ctD(X,\blank),A) \simeq AX\) del lemma di Yoneda.
\end{proposition}

\begin{proof}
 L'idea della dimostrazione è che dato un cocono \(\{\bar g_{(X,x)} \colon \ctD(X,\blank) \nat G\}_{(X,x) \in \Elts\ctD A}\), la sua fattorizzazione attraverso il cocono \(\{\xi \colon \ctD(X,\blank) \nat A\}_{(X,x) \in \Elts\ctD A}\) è data dalla trasformazione naturale \(g \colon A \nat G\) la cui componente in \(X \in \ctD\) è definita da \(g_X(x) = \bar{g}_{(X,x)}\), dove \(\bar{g}_{(X,x)} \in GX\) è l'elemento che corrisponde a \(g_{(X,x)}\) nella biiezione naturale \(\Nat(\ctD(X,\blank),G) \natIso GX\) del lemma di Yoneda.
 \Todo{...}
\end{proof}

\begin{remark}\label{oss_colim_rappr}
 \hfill
 \begin{enumerate}
  \item Poiché questo ci servirà nell'esempio \ref{exem_proreg}, osserviamo che, combinando la proposizione \ref{prop_colim_rappr}
        e il teorema \ref{th_colim_piccoli}, otteniamo che \(A \colon \ctD \fun \ctSet\) è un quoziente regolare di un coprodotto di funtori rappresentabili
        \[\xymatrix{p_A \colon \cop_{(X,x) \in \Elts\ctD A}\ctD(X,\blank) \ar[r] & A.}\]
        La somma a dominio esiste, perché la famiglia su cui viene indicizzata è piccola, dato che lo è $\Elts\ctD A$ (la classe dei suoi oggetti ha cardinalità al più $|\ctD_0|\cdot\sup_X |AX|$).
  \item Osserviamo anche che, poiché (come osservato in \ref{}) la categoria degli elementi \(\Elts\ctD A\) e la comma \(\coyon_{\ctD}/A\) sono isomorfe, avremmo potuto evitare di introdurre \(\Elts\ctD A\) ed esprimere tutto usando solo la comma \(\coyon_{\ctD}/A\). Usare la categoria degli elementi però rende un po' più semplice la descrizione del funtore dimenticante \(\Phi_A\) (che sarebbe, comunque, nient'altro che una delle proiezioni dalla comma).
 \end{enumerate}
\end{remark}

Il fatto che ogni funtore \(A \colon \ctD \fun \ctSet\) sia un colimite di funtori rappresentabili ci permette ora di dimostrare che l'immersione
di Yoneda \(\coyon_{\ctD} \colon \ctD^{\op} \fun \ctSet^{\ctD}\) è il completamento di \(\ctD^{\op}\) per colimiti, nel senso della seguente proprietà universale.

\begin{proposition}\label{prop_univ_yon}
 Siano \(\ctD\) una categoria piccola e \(\ctB\) una categoria piccolo-cocompleta nel senso di \ref{}. Per ogni funtore \(F \colon \ctD^{\op} \fun \ctB\) esiste un funtore
 \(F^* \colon \ctSet^{\ctD} \fun \ctB\) che preserva i colimiti e tale che il diagramma
 \[\xymatrix{\ctD^{\op} \ar[r]^-{\coyon_{\ctD}} \ar[d]_{F} & \ctSet^{\ctD} \ar[ld]^{F^*} \\
  \ctB}\]
 commuta a meno di un isomorfismo naturale. Inoltre un tale funtore \(F^*\) è essenzialmente unico.
\end{proposition}

\begin{proof}
 Poiché \(F^*\) deve preservare i colimiti e \(F^* \cmp \coyon_{\ctD} \simeq F\), la proposizione \ref{prop_colim_rappr}
 fornisce in modo necessario la definizione di \(F^*\) sugli oggetti di \(\ctSet^{\ctD}\)
 \[F^*(A) = F^*(\colim_{\Elts\ctD A}(\coyon_{\ctD} \cmp \Phi_A)) \simeq \colim_{\Elts\ctD A}(F^* \cmp \coyon_{\ctD} \cmp \Phi_A))
  \simeq \colim_{\Elts\ctD A}(F \cmp \Phi_A)\]
 Per quanto riguarda la definizione di \(F^*\) sulle frecce di \(\ctSet^{\ctD}\), una trasformazione naturale \(\lambda \colon A \nat A'\)
 induce un funtore
 \[\Phi_{\lambda} \colon \Elts\ctD A \fun \mathrm{El}(A'), \;
  \Phi_{\lambda}(f \colon (X,x) \to (Y,y)) = (f \colon (X,\lambda_X(x)) \to (Y,\lambda_Y(y))\]
 che fa commutare il diagramma
 \[\xymatrix{\Elts\ctD A \ar[rr]^-{\Phi_{\lambda}} \ar[rd]_{\Phi_A} & & \mathrm{El}(A') \ar[ld]^{\Phi_{A'}} \\
  & \ctD^{\op}}\]
 La proprietà universale del colimite \(F^*(A)\) fornisce ora un'unica freccia \(F^*(\lambda)\) che fa commutare, per ogni
 \((X,x) \in \Elts\ctD A\), il seguente diagramma (dove le frecce \(\sigma\) sono le iniezioni nei colimiti)
 \[\xymatrix{F^*(A) \ar[rr]^-{F^*(\lambda)} & & F^*(A') \\
  & F(\Phi_A(X,x)) \ar[lu]^{\sigma^A_{(X,x)}} \ar[ru]_{\sigma^{A'}_{\Phi_{\lambda}(X,x)}} }\]
 Verifichiamo ora l'isomorfismo \(F^* \cmp \ctY_{\ctD} \simeq F\). Se \(A = \ctD(Z,\blank)\) per un oggetto \(Z \in \ctD\), oggetti e
 frecce della categoria \(\Elts\ctD A\) assumono la seguente forma
 \[f \colon (X,x \colon Z \to X) \to (Y, y \colon Z \to Y) \;\mbox{ tale che }\; f \cmp y = x \]
 In particolare, \((Z,\id \colon Z \to Z)\) è l'oggetto terminale di \(\Elts\ctD A\) poiché, per ogni altro oggetto \((X,x)\),
 si ha un'unica freccia data da \(x \colon (X,x) \to (Z,\id)\). Ora possiamo mostrare che
 \[ \colim_{\Elts\ctD A}(F \cmp \Phi_A) = \langle FZ,\{F(x) \colon FX \to FZ\}_{(X,x) \in \Elts\ctD A} \rangle \]
 La naturalità di tale famiglia segue subito dalla condizione sulle frecce di \(\Elts\ctD A\). Per quanto riguarda l'universalità,
 se \(\{b_{(X,x)} \colon FX \to B\}_{(X,x) \in \Elts\ctD A}\) è un'altra famiglia naturale, si ha che
 \[ \xymatrix{FZ \ar[rr]^-{b_{(Z,\id)}} & & B \\ & FX \ar[lu]^{F(x)} \ar[ru]_{b_{(X,x)}} } \]
 commuta, il che indica che \(b_{(Z,\id)} \colon FZ \to B\) è una fattorizzazione della famiglie \(\{b_{(X,x)}\}\) attraverso la famiglia
 \(\{F(x)\}\). Tale fattorizzazione è unica: se \(t \colon FZ \to B\) è tale che \(t \cmp F(x) = b_{(X,x)}\) per ogni
 \((X,x) \in \Elts\ctD A\), in particolare \(t \cmp F(\id) = b_{(Z,\id)}\) e quindi \(t = b_{(Z,\id)}\). \\
 Rimane da verificare che \(F^*\) preserva i colimiti, ma questo segue dal fatto che \(F^*\) ammette un aggiunto destro
 \(R_F \colon \ctB \fun \ctSet^{\ctD}\) definito da
 \[ R_F(B) = \ctB(F-,B) \colon \ctD \fun \ctSet \]
 La biiezione naturale che garantisce l'aggiunzione \(F^* \dashv R_F\) si ottiene così:
 \[ \ctB(F^*A,B) = \ctB(\colim_{\Elts\ctD A}(F \cmp \Phi_A),B) \simeq \lim_{\Elts\ctD A}\ctB(FX,B) \simeq \]
 \[ \simeq \lim_{\Elts\ctD A} \Nat(\ctD(X,\blank), \ctB(F-,B)) \simeq \Nat(\colim_{\Elts\ctD A}(\ctY_{\ctD} \cmp \Phi_A), \ctB(F-,B))
  = \Nat(A,R_F(B)) \]
 avendo usato, per passare dalla prima alla seconda riga, le biiezione del lemma di Yoneda.
\end{proof}

Il prossimo obiettivo è di utilizzare la proprietà universale \ref{prop_univ_yon} per caratterizzare le categorie del tipo
\(\ctSet^{\ctD}\) (proposizione \ref{prop_caratt_pref}). Ci occorre una definizione preliminare in cui introduciamo gli oggetti
assolutamente, perfettamente e finitamente presentabili. Solo i primi interverranno nella caratterizzazione \ref{prop_caratt_pref},
gli altri saranno utilizzati nel capitolo \ref{cap_cat_alg}.

\begin{definition}\label{def_finpres_perfpres}
 Consideriamo una categoria \(\ctA\), un oggetto \(A \in \ctA\) e il funtore rappresentabile \(\ctA(A,\blank) \colon \ctA \fun \ctSet\).
 L'oggetto \(A\) è
 \begin{enumerate}
  \item assolutamente presentabile se il funtore \(\ctA(A,\blank)\) preserva i colimiti,
  \item perfettamente presentabile se il funtore \(\ctA(A,\blank)\) preserva i colimiti setacciati,
  \item finitamente presentabile se il funtore \(\ctA(A,\blank)\) preserva i colimiti filtrati.
 \end{enumerate}
\end{definition}

Ovviamente un oggetto assolutamente presentabile è perfettamente presentabile e un oggetto perfettamente presentabile è
finitamente presentabile. Nelle due prossime proposizioni vediamo qualche proprietà di stabilità degli oggetti assolutamente,
perfettamente e finitamente presentabili.

\begin{proposition}\label{prop_stab_present}
 Sia \(\ctA\) una categoria.
 \begin{enumerate}
  \item Se \(A \in \ctA\) è un oggetto assolutamente presentabile e se \(R\) è un retratto di \(A\), allora anche \(R\) è assolutamente
        presentabile. Lo stesso vale per gli oggetti perfettamente presentabili e per gli oggetti finitamente presentabili.
  \item Un colimite finito di oggetti finitamente presentabili è finitamente presentabile.
  \item Un coprodotto finito di oggetti perfettamente presentabili è perfettamente presentabile.
 \end{enumerate}
\end{proposition}

\begin{proof}
 1. Se \(R\) è un retratto di \(A\), allora il funtore \(\ctA(R,\blank)\) è un retratto del funtore \(\ctA(A,\blank)\). La proposizione
 \ref{prop_funt_retratto} permette di concludere. \\
 2. Sia \(I \colon \ctI \fun \ctA\) un funtore definito su una categoria \(\ctI\) finita e sia \(L\) il colimite di \(I\). Supponiamo che,
 per ogni \(i \in \ctI\), l'oggetto \(I(i)\) sia finitamente presentabile. Consideriamo anche un funtore \(F \colon \ctD \fun \ctA\),
 dove \(\ctD\) è una categoria filtrata. Il fatto che \(L\) sia finitamente presentabile si mostra con la seguente catena di
 isomorfismi naturali:
 \[\begin{array}{rclcl}
   \ctA(L,\colim_{\ctD}FD) & \simeq & \ctA(\colim_{\ctI}Ii,\colim_{\ctD}FD) &  & (\mbox{definizione di } L)                \\
                           & \simeq & \lim_{\ctI}\ctA(Ii,\colim_{\ctD}FD)   &  & (\mbox{osservazione \ref{oss_rappr_lim}}) \\
                           & \simeq & \lim_{\ctI}(\colim_{\ctD}\ctA(Ii,FD)) &  & (Ii \mbox{ è finitamente presentabile})   \\
                           & \simeq & \colim_{\ctD}(\lim_{\ctI}\ctA(Ii,FD)) &  & (\mbox{teorema \ref{th_caratt_filtr}})    \\
                           & \simeq & \colim_{\ctD}\ctA(\colim_{\ctI}Ii,FD) &  & (\mbox{osservazione \ref{oss_rappr_lim}}) \\
                           & \simeq & \colim_{\ctD}\ctA(L,FD)               &  & (\mbox{definizione di } L)
  \end{array}\]
 3. Simile alla prova del punto 3 prendendo come \(\ctI\) una categoria finita e discreta, come \(\ctD\) una categoria setacciata
 e usando il teorema \ref{th_caratt_sifted}.
\end{proof}

\begin{proposition}\label{prop_stab_perf_pres}
 Consideriamo un'aggiunzione
 \[ \xymatrix{\ctB \ar@<-0.5ex>[r]_-{G} & \ctA \ar@<-0.5ex>[l]_-{F} } \;\;\; F \dashv G \]
 Se il funtore \(G\) preserva i colimiti (rispettivamente, i colimiti setacciati, i colimiti filtrati) e se \(A \in \ctA\) è assolutamente
 presentabile (rispettivamente, perfettamente presentabile, finitamente presentabile), allora anche \(FA \in \ctB\) lo è.
\end{proposition}

\begin{proof}
 Vediamo il caso degli oggetti assolutamente presentabili. Il diagramma
 \[ \xymatrix{\ctB \ar[rr]^-{G} \ar[rd]_{\ctB(FA,\blank)} & & \ctA \ar[ld]^{\ctA(A,\blank)} \\ & \ctSet} \]
 commuta a meno dell'isomorfismo naturale fornito dall'aggiunzione \(F \dashv G\). Ne segue che \(\ctB(FA,\blank)\) preserva i
 colimiti in quanto composto di due funtori che li preservano.
\end{proof}

\begin{example}\label{ex_ass_present_pref}
 In questo esempio vogliamo identificare gli oggetti assolutamente presentabili della categoria \(\ctSet^{\ctC}\), dove
 \(\ctC\) è una categoria piccola.
 Come primo passo, osserviamo che i funtori rappresentabili sono degli oggetti assolutamente presentabili di
 \(\ctSet^{\ctC}\). Infatti, nella situazione
 \[\xymatrix{\ctD \ar[r]^-{H} & \ctSet^{\ctC} \ar[rr]^-{\Nat(\ctC(X,\blank),\blank)} & & \ctSet}\]
 usando due volte il lemma di Yoneda e il fatto che i colimiti in \(\ctSet^{\ctC}\) si calcolano punto a
 punto (proposizione \ref{prop_lim_puntuali}), si ha
 \[\Nat(\ctC(X,\blank),\colim_{\ctD}H) \simeq (\colim_{\ctD}H)(X) = \colim_{\ctD}((HD)(X)) \simeq \colim_{\ctD}\Nat(\ctC(X,\blank),HD))\]
 il che conclude la prova che i rappesentabili son assolutamente presentabili.
 La proposizione \ref{prop_stab_present} implica ora che i retratti dei funtori rappresentabili sono oggetti
 assolutamente presentabili di \(\ctSet^{\ctC}\). Non ce ne sono altri: se \(A \in \ctSet^{\ctC}\) è assolutamente
 presentabile, allora \(A\) è un retratto di un funtore rappresentabile. Infatti, con la notazione della proposizione
 \ref{prop_colim_rappr}, si ha
 \[ \Nat(A,A) = \Nat(A, \colim_{\Elts\ctD A}(\ctY_{\ctC} \cmp \Phi_A)) \simeq \colim_{\Elts\ctD A} \Nat(A, \ctC(X,\blank)) \]
 Poiché \(\id_A \in \Nat(A,A)\), tale biiezione implica in particolare che esistono \((X,x) \in \Elts\ctD A\) e
 \(\alpha \in \Nat(A,\ctC(X,\blank))\) tali che \(\xi \cmp \alpha = \id_A\), dove \(\xi \colon \ctC(X,\blank) \nat A\) è la trasformazione
 naturale che corrisponde all'elemento \(x \in A(X)\). In conclusione, \(A\) è un retratto di \(\ctC(X,\blank)\).
\end{example}

\begin{lemma}\label{lemma_caratt_pref}
 Nella situazione della proposizione \ref{prop_univ_yon}
 \[\xymatrix{\ctD^{\op} \ar[r]^-{\ctY_{\ctD}} \ar[d]_{F} & \ctSet^{\ctD} \ar[ld]^{F^*} \\
  \ctB}\]
 se \(F\) è pieno e fedele e se, per ogni \(Z \in \ctD\), l'oggetto \(FZ\) è assolutamente presentabile, allora \(F^*\) è pieno e fedele.
\end{lemma}

\begin{proof}
 Cominciamo con l'osservare cha la proposizione \ref{prop_colim_rappr} permette di ridurre lo studio del carattere pieno e fedele
 di \(F^*\) alle frecce il cui dominio è un funtore rappresentabile. Infatti, per \(A,A' \in \ctSet^{\ctD}\) si ha
 \[ \Nat(A,A') = \Nat(\colim_{\Elts\ctD A}(\ctY_{\ctD} \cmp \Phi_A),A') \simeq \lim_{\Elts\ctD A}\Nat(\ctD(X,\blank),A') \]
 Fissiamo dunque \(A \in \ctSet^{\ctD}\) e \(Z \in \ctD\) e consideriamo il, seguente diagramma, la cui commutatività per ogni
 \((X,x) \in \Elts\ctD A\) definisce l'azione di \(F^*\) sulle frecce
 \[\xymatrix{ \Nat(\ctD(Z,\blank), A) \ar[rr]^-{F^*} \ar@{=}[d] & & \ctB(FZ, F^*A) \ar@{=}[d] \\
  \Nat(\ctD(Z,\blank), \colim_{\Elts\ctD A}(\ctY_{\ctD} \cmp \Phi_A)) & & \ctB(FZ, \colim_{\Elts\ctD A}(F \cmp \Phi_A)) \\
  \colim_{\Elts\ctD A}\Nat(\ctD(Z,\blank),\ctD(X,\blank)) \ar[u]^{\mathrm{can}} & & \colim_{\Elts\ctD A}\ctB(FZ,FX) \ar[u]_{\mathrm{can}} \\
  \Nat(\ctD(Z,\blank)\ctD(X,\blank)) \simeq \ctD(Z,X) \ar[rr]_-{F} \ar[u]^{\sigma_{(X,x)}} & & \ctB(FZ,FX) \ar[u]_{\sigma_{(X,x)}} }\]
 Le frecce chiamate \(\sigma\) sono come al solito le iniezioni nei colimiti e l'isomorfismo è la biiezione del lemma di Yoneda.
 Osserviamo ora che le due frecce canoniche sono degli isomorfismi: quella di sinistra perché \(\ctD(Z,\blank)\) è assolutamente presentabile
 (esempio \ref{ex_ass_present_pref}) e quella di destra perché \(FZ\) è assolutamente presentabile per ipotesi. Poiché la freccia alla base
 del diagramma è un isomorfismo (perché \(F\) è pieno e fedele), anche la feccia in alto è un isomorfismo, cioé \(F^*\) è pieno e fedele.
\end{proof}

\begin{proposition}\label{prop_caratt_pref}
 Una categoria \(\ctB\) è equivalente a una categoria del tipo \(\ctSet^{\ctD}\) per \(\ctD\) piccola se e solo se \(\ctB\) è cocompleta e contiene
 un insieme \(\ctG\) di oggetti assolutamente presentabili tale che ogni oggetto di \(\ctB\) è colimite di oggetti di \(\ctG\).
\end{proposition}

\begin{proof}
 Sappiamo che \(\ctSet^{\ctD}\) è cocompleta (corollario \ref{cor_lim_cat_pref}) e come insieme \(\ctG\) si può prendere
 l'insieme dei funtori \(\ctD \fun \ctSet\) rappresentabili, vedi proposizione \ref{prop_colim_rappr}. \\
 Viceversa,  consideriamo come categoria \(\ctD\) la duale della sottocategoria piena di \(\ctB\) che ha come oggetti quelli
 dell'insieme \(\ctG\) e prendiamo come funtore \(F \colon \ctD^{\op} \fun \ctB\) l'inclusione piena. Si può applicare il lemma
 \ref{lemma_caratt_pref}, per cui il funtore indotto \(F^* \colon \ctSet^{\ctD} \fun \ctB\) è pieno e fedele. Inoltre, per ogni
 oggetto \(B \in \ctB\) esistono per ipotesi una categoria piccola \(\ctD_B\) e un funtore \(\Psi_B \colon \ctD_B \fun \ctD^{\op}\) tali che
 \[ B = \colim_{\ctD_B}(F \cmp \Psi_B) \]
 Se poniamo \(A = \colim_{\ctD_B}(\ctY_{\ctD} \cmp \Psi_B)\) e usiamo il fatto che \(F^*\) preserva i colimiti, otteniamo
 \[ F^*A = F^*(\colim_{\ctD_B}(\ctY_{\ctD} \cmp \Psi_B)) \simeq \colim_{\ctD_B}(F^* \cmp \ctY_{\ctD} \cmp \Psi_B) \simeq
  \colim_{\ctD_B}(F \cmp \Psi_B) = B \]
 Questo mostra che \(F^* \colon \ctSet^{\ctD} \fun \ctB\) è essenzialmente suriettivo ed è quindi un'equivalenza.
\end{proof}

Terminiamo questa sezione con una proprietà di stabilità per le categorie di prefasci.

\begin{corollary}\label{cor_stab_pref}
 Sia \(\ctC\) una categoria piccola e sia \(i \colon \ctA \fun \ctSet^{\ctC}\) una sottocategoria riflessiva e chiusa per colimiti.
 La categoria \(\ctA\) è a sua volta equivalente a una categoria di prefasci su una categoria piccola.
\end{corollary}

\begin{proof}
 Mostriamo che \(\ctA\) soddisfa le ipotesi della caratterizzazione \ref{prop_caratt_pref}. Certamente \(\ctA\) è cocompleta
 in quanto sottocategoria riflessiva di una categoria cocompleta. Sia ora \(R \colon \ctSet^{\ctC} \fun \ctA\) l'aggiunto sinistro.
 e considero l'insieme di oggetti di \(\ctA\) definito da
 \[\ctG = \{R(\ctC(X,\blank)) \mid X \in \ctC\}\]
 In virtù della proposizione \ref{prop_stab_perf_pres} e dell'esempio \ref{ex_ass_present_pref}, tali oggetti sono assolutamente
 presentabili in \(\ctA\). Per terminare, consideriamo un oggetto \(A \in \ctA\). Per la proposizione \ref{prop_colim_rappr},
 sappiamo che \(iA = \colim_{\Elts\ctD A}(\ctY_{\ctC} \cmp \Phi_A)\). Usando che \(i\) è pieno e fedele (e quindi la
 counità \(R(iA) \to A\) dell'aggiunzione è un isomorfismo) e che \(R\) preserva i colimiti, otteniamo
 \[ A \simeq R(iA) = R(\colim_{\Elts\ctD A}(\ctY_{\ctC} \cmp \Phi_A))
  \simeq \colim_{\Elts\ctD A}(R \cmp \ctY_{\ctC} \cmp \Phi_A) \]
 il che conclude la prova.
\end{proof}

\begin{exercise}\label{ex_perf_present_pref}
 Gli oggetti perfettamente presentabili in una categoria della forma \(\ctSet^{\ctC}\) si possono caratterizzare in modo analogo a
 quanto fatto in \ref{ex_ass_present_pref} con gli oggetti assolutamente presentabili: i prefasci perfettamente presentabili sono esattamente i retratti dei coprodotti finiti di rappresentabili. 
 
 Vediamo le tappe essenziali della prova e lasciamo a chi legge il compito di completarla. 
 \begin{itemize}
      \item I prefasci rappresentabili sono assolutamente presentabili e quindi perfettamente presentabili.
 Inoltre, i coprodotti finiti di oggetti perfettamente presentabili sono perfettamente presentabili. Infine, i retratti di oggetti
 perfettamente presentabili sono perfettamente presentabili.
 \item Viceversa, per ogni \(A \in \ctSet^{\ctC}\) esiste una categoria piccola \(\ctD_A\) e un funtore \(\Psi_A \colon \ctD_A \fun \ctC^{\op}\)
 tale che \(A = \colim_{\ctD_A}(\ctY_{\ctC} \cmp \Psi_A)\). Ne segue che \(A\) è il coequalizzatore di un grafo riflessivo che ha come
 codominio il coprodotto \(\cop \ctC(\Psi_AD,\blank)\) indicizzato dagli oggetti \(D \in \ctD_A\). Tale coprodotto è anche il
 colimite filtrante dei suoi sottocoprodotti finiti (proposizione \ref{prop_coprod_colimfiltr})
 \[ \cop_{D \in \ctD_A}\ctC(\psi_AD,\blank) =  \colim_{I \in \ctP_{\mathrm{f}}(\ctD_A)} \left(\cop_{D \in I}\ctC(\Psi_AD,\blank)\right) \]
 Se ora assumiamo che \(A\) è perfettamente presentabile e applichiamo il funtore \(\Nat(A,\blank)\), otteniamo una funzione suriettiva
 \[ \colim_{I \in \ctP_{\mathrm{f}}(\ctD_A)}\Nat\left(A, \cop_{D \in I}\ctC(\Psi_AD,\blank)\right) \to \Nat(A,A) \]
 Poiché \(\id \in \Nat(A,A)\), ne segue che esistono \(I \in \ctP_{\mathrm{f}}(\ctD_A)\) e
 \(\lambda \in \Nat\left(A,\cop_{D \in I}\ctC(\Psi_AD,\blank)\right)\) tali che il seguente diagramma commuta
  \[\xymatrix{ 
   A\ar@{=}[r]\ar[d]_\lambda & A \\    
   \cop_{D \in I}\ctC(\Psi_AD,\blank)\ar[r] & \cop_{D \in \ctD_A}\ctC(\Psi_AD,\blank)\ar[u]
  }\]
 il che mostra che \(A\) è un  retratto di un coprodotto finito di prefasci rappresentabili.
 \end{itemize}
\end{exercise}

\section{Equivalenze fra categorie di prefasci}\label{sec_equiv_prefasci}

Vista l'importanza delle categorie del tipo \(\ctSet^{\ctC}\) con \(\ctC\) piccola, ci si può chiedere quando due categorie
piccole \(\ctC\) e \(\ctD\) danno luogo a categorie di prefasci equivalenti: \(\ctSet^{\ctC} \simeq \ctSet^{\ctD}\).
Per rispondere in modo esaustivo a tale domanda dobbiamo rivisitare il completamento per idempotenti (definizione
\ref{def_idemp_compl} e proposizione \ref{prop_compl_idemp}) nel caso di una categoria piccola.

\begin{corollary}\label{cor_idcompl_ap}
 Sia \(\ctC\) una categoria piccola.
 \begin{enumerate}
  \item Il completamento per idempotenti \(\mathrm{Ic}(\ctC)\) è una categoria piccola.
  \item Il completamento per idempotenti di \(\ctC^{\op}\) è equivalente alla corestrizione dell'immersione di Yoneda \(\ctY_\ctC\)
        alla sottocategoria \(\ctSet^{\ctC}_{\mathrm{ap}}\) di \(\ctSet^{\ctC}\) degli oggetti assolutamente presentabili.
  \item Se \(\ctC\) è completa per idempotenti, la corestrizione
        \(\ctY_{\ctC} \colon \ctC^{\op} \fun \ctSet^{\ctC}_{\mathrm{ap}}\) è un'equivalenza.
 \end{enumerate}
\end{corollary}

\begin{proof}
 1. \(\mathrm{Ic}(\ctC)\) è localmente piccola perché \(\mathrm{Ic}(\ctC)((X \to X),(Y \to Y)) \subseteq \ctC(X,Y)\).
 Inoltre, gli oggetti di \(\ctC\) costituiscono un insieme e, per ogni oggetto \(X\), le frecce idempotenti \(X \to X\)
 costituiscono anch'esse un insieme. Quindi \(\mathrm{Ic}(\ctC)\) ha un insieme di oggetti. \\
 2. Un retratto di un oggetto assolutamente presentabile è assolutamente presentabile (proposizione \ref{prop_stab_present}.
 Ne segue che \(\ctSet^{\ctC}_{\mathrm{ap}}\) è chiusa in \(\ctSet^{\ctC}\) per retratti ed è quindi completa per idempotenti
 (osservazione \ref{oss_idCompl_retr}). Ora possiamo utilizzare la proprietè universale \ref{prop_compl_idemp} per completare
 il diagramma
 \[ \xymatrix{\ctC^{\op} \ar[rr]^-{\mathrm{Ic}_{\ctC^{\op}}} \ar[rd]_{\ctY_{\ctC}} & & \mathrm{Ic}(\ctC^{\op}) \ar[ld]^{\ctY_{\ctC}'} \\
  & \ctSet^{\ctC}_{\mathrm{ap}} } \]
 Grazie all'esempio \ref{ex_ass_present_pref} possiamo applicare il primo punto del lemma \ref{lemma_caratt_idcompl} e
 otteniamo che \(\ctY_{\ctC}' \colon \mathrm{Ic}(\ctC^{\op}) \to \ctSet^{\ctC}_{\mathrm{ap}}\) è un'equivalenza. \\
 3. Se \(\ctC\) è completa per idempotenti, nel diagramma precedente anche \(\mathrm{Ic}_{\ctC^{\op}}\) è un'equivalenza.
\end{proof}

Possiamo ora formulare una prima risposta alla domanda che ci siamo posti all'inizio di questa sezione.

\begin{proposition}\label{prop_morita_prefasci}
 Siano \(\ctC\) e \(\ctD\) due categorie piccole.
 \begin{enumerate}
  \item Il funtore \(\mathrm{Ic}_{\ctC} \colon \ctC \fun \mathrm{Ic}(\ctC)\) induce un'equivalenza
        \(\ctSet^{\mathrm{Ic}_{\ctC}} \colon \ctSet^{\mathrm{Ic}(\ctC)} \fun \ctSet^{\ctC}\).
  \item  \(\ctSet^{\ctC} \simeq \ctSet^{\ctD}\) se e solo se \(\ctSet^{\ctC}_{\mathrm{ap}} \simeq \ctSet^{\ctD}_{\mathrm{ap}}\)
        se e solo se \(\mathrm{Ic}(\ctC) \simeq \mathrm{Ic}(\ctD)\).
 \end{enumerate}
\end{proposition}

\begin{proof}
 1. Poiché \(\ctSet\) è completa e quindi completa per idempotenti, questo punto è un  caso particolare
 del corollario \ref{cor_equiv_idcompl}. \\
 2. L'equivalenza fra le condizioni \(\ctSet^{\ctC}_{\mathrm{ap}} \simeq \ctSet^{\ctD}_{\mathrm{ap}}\) e
 \(\mathrm{Ic}(\ctC) \simeq \mathrm{Ic}(\ctD)\) segue subito dall'osservazione \ref{oss_idcompl_duale} e dal
 corollario \ref{cor_idcompl_ap}. \\
 Se \(\mathrm{Ic}(\ctC) \simeq \mathrm{Ic}(\ctD)\), allora usando il primo punto si ha
 \(\ctSet^{\ctC} \simeq \ctSet^{\mathrm{Ic}(\ctC)} \simeq \ctSet^{\mathrm{Ic}(\ctD)} \simeq \ctSet^{\ctD}\). \\
 Viceversa, per la proposizione \ref{prop_stab_perf_pres} ogni equivalenza \(\ctSet^{\ctC} \simeq \ctSet^{\ctD}\)
 si restringe a un'equivalenza \(\ctSet^{\ctC}_{\mathrm{ap}} \simeq \ctSet^{\ctD}_{\mathrm{ap}}\) fra le sottocategorie
 degli oggetti assolutamente presentabili.
\end{proof}

\begin{remark}\label{oss_dualita}
 La proposizione \ref{prop_morita_prefasci} ci dice in particolare che, se due categorie piccole \(\ctC\) e \(\ctD\) sono complete
 per idempotenti,allora \(\ctSet^\ctC \simeq \ctSet^\ctD\) se e solo se \(\ctC \simeq \ctD\). Questo fatto si può anche inquadrare
 in una dualità fra la 2-categoria delle categorie piccole e complete per idempotenti e la 2-categoria delle categorie di prefasci.
 Senza entrare in una spiegazione dettagliata di tale risultato, che richiede evidentemente di introdurre le 2-categorie e i 2-funtori,
 ci limitiamo a enunciare e dimostrare i due ingredienti essenziali di tale dualità, ovvero le proposizioni \ref{prop_funt_indotto_pref}
 e \ref{prop_dualita_cat_pref}.
\end{remark}

\begin{proposition}\label{prop_funt_indotto_pref}
 Ogni funtore \(M \colon \ctC \fun \ctD\) fra categorie piccole induce un funtore
 \[ \ctSet^M \colon \ctSet^{\ctD} \fun \ctSet^{\ctC}, \; \ctSet^M(B) = B \cmp M \colon \ctC \fun \ctD \fun \ctSet \]
 che preserva i colimiti e ha un aggiunto sinistro che denoteremo \(M^* \colon \ctSet^{\ctC} \fun \ctSet^{\ctD}\).
\end{proposition}

\begin{proof}
 Il fatto che \(\ctSet^M\) preservi i colimiti è un caso particolare del corollario \ref{cor_funt_indotto_catfunt}.
 Per mostrare che \(\ctSet^M\) ha aggiunto sinistro, consideriamo il seguente diagramma
 \[ \xymatrix{\ctC^{\op} \ar[d]_{M^{\op}} \ar[r]^-{\ctY_{\ctC}} \ar@{.>}[rd]^{F} & \ctSet^{\ctC} \ar[d]^{M^*} \\
  \ctD^{\op} \ar[r]_-{\ctY_{\ctD}} & \ctSet^{\ctD}} \]
 dove abbiamo applicato la proprietà universale dell'immersione di Yoneda \(\ctY_{\ctC}\) (proposizione
 \ref{prop_univ_yon}) al funtore \(F = \ctY_{\ctD} \cmp M^{\op}\) per ottenere il funtore \(M^*\) che preserva
 i colimiti e che rende il diagramma commutativo a meno di isomorfismo naturale. Sappiamo che \(M^*\) ha
 un aggiunto destro \[ R_M \colon \ctSet^{\ctD} \fun \ctSet^{\ctC}, \; R_M(B) = \Nat(F-,B) \colon \ctC \fun \ctSet \]
 e resta da mostrare che \(R_M\) è isomorfo a \(\ctSet^M\). Per ogni \(X \in \ctC\) si ha
 \[ (R_M(B))(X) = \Nat(FX,B) = \Nat(\ctD(MX,\blank),B) \simeq B(MX) = (\ctSet^M(B))(X) \]
 avendo usato alla penultima tappa il lemma di Yoneda.
\end{proof}

\begin{example}\label{ex_M_subcat_refl}
 Esaminiamo la dimostrazione della proprietà di stabilità delle categorie di prefasci (corollario \ref{cor_stab_pref})
 dal punto di vista dei funtori di tipo \(\ctSet^M\). Se \(i \colon \ctA \fun \ctSet^{\ctC}\) è una sottocategoria riflessiva,
 con riflettore \(R\), abbiamo posto \(\ctD^{\op} = \{R(\ctC(X,\blank)) \mid X \in \ctC\}\). Se ora chiamiamo \(F\)
 l'inclusione piena \(\ctD \subseteq \ctA\), abbiamo mostrato che l'unico funtore \(F^* \colon \ctSet^{\ctD} \fun \ctA\)
 che preserva i colimiti e tale che \(F^* \cmp \ctY_{\ctD} \simeq F\) è un'equivalenza. Se definiamo
 \(M \colon \ctC \fun \ctD\) come la corestrizione di \(R \cmp \ctY_{\ctD}\) a \(\ctD\), abbiamo che \(F^* \cmp M^* \simeq R\).
 Infatti
 \[F^* \cmp M^* \cmp \ctY_{\ctC} \simeq F^* \cmp \ctY_{\ctD} \cmp M^{\op} \simeq F \cmp M^{\op} = R \cmp \ctY_{\ctC}\]
 e tanto \(F^* \cmp M^*\) che \(R\) preservano i colimiti. Possiamo concludere che, a meno dell'equivalenza
 \(\ctA \simeq \ctSet^{\ctD}\), l'aggiunzione di partenza \(R \dashv i\) coincide con l'aggiunzione indotta \(M^* \dashv \ctSet^M\).
\end{example}

\begin{proposition}\label{prop_dualita_cat_pref}
 Sia \(M \colon \ctC \fun \ctD\) un funtore fra categorie piccole.
 Se \(\ctC\) e \(\ctD\) sono complete per idempotenti, allora la costruzione
 \[ M \colon \ctC \fun \ctD \; \mapsto \; \ctSet^M \colon \ctSet^{\ctD} \fun \ctSet^{\ctC} \]
 della proposizione \ref{prop_funt_indotto_pref} realizza un'equivalenza fra la categoria dei funtori da
 \(\ctC\) a \(\ctD\) e la categoria dei funtori da \(\ctSet^{\ctD}\) a
 \(\ctSet^{\ctC}\) che preservano i colimiti e hanno aggiunto sinistro.
\end{proposition}

\begin{proof}
 Data un'aggiunzione
 \[ \xymatrix{\ctSet^{\ctD} \ar@<-0.5ex>[r]_-{R} & \ctSet^{\ctC} \ar@<-0.5ex>[l]_-{L} } \;\;\; L \dashv R \]
 assumiamo che \(R\) preservi i colimiti e cerchiamo un funtore \(M \colon \ctC \fun \ctD\) tale che \(R \simeq \ctSet^M\).
 Per questo, consideriamo il diagramma
 \[ \xymatrix{\ctC^{\op} \ar@{.>}[ddd]_{M^{\op}} \ar[rr]^-{\ctY_{\ctC}} \ar[rd]^{\simeq} & & \ctSet^{\ctC} \ar[ddd]^{L} \\
  & \ctSet^{\ctC}_{\mathrm{ap}} \ar@{.>}[d]^{L_{\mathrm{ap}}} \ar[ru] \\
  & \ctSet^{\ctD}_{\mathrm{ap}} \ar[rd] \\
  \ctD^{op} \ar[ru]_{\simeq} \ar[rr]_-{\ctY_{\ctD}} & & \ctSet^{\ctD} } \]
 Poiché \(\ctC\) e \(\ctD\) sono complete per idempotenti, le corestrizioni di \(\ctY_{\ctC}\) e di \(\ctY_{\ctD}\) alle
 sottocategorie dei prefasci assolutamente presentabili sono delle equivalenze (corollario \ref{cor_idcompl_ap}).
 Inoltre, per la proprietà \ref{prop_stab_perf_pres}, l'aggiunto sinistro \(L\) si restringe alle sottocategorie dei prefasci
 assolutamente presentabili. Chiamiamo \(L_{\mathrm{ap}}\) tale restrizione e poniamo
 \[ M^{\op} \colon \ctC^{\op} \simeq \ctSet^{\ctC}_{\mathrm{ap}} \xto{L_{\mathrm{ap}}} \ctSet^{\ctD}_{\mathrm{ap}} \simeq \ctD^{\op} \]
 Poiché il diagramma completo commuta, il funtore \(L\) è l'estensione di \(\ctY_{\ctD} \cmp M^{\op}\) fornita dalla
 proprietà universale dell'immersione di Yoneda \(\ctY_{\ctC}\) (proposizione \ref{prop_univ_yon}) e quindi il suo aggiunto destro
 \(R\) coincide, a meno di iso naturale, con \(\ctSet^M\), come mostrato nella prova della proposizione \ref{prop_funt_indotto_pref}.
\end{proof}

Dal corollario \ref{cor_equiv_sutr} sappiamo che, date tre categorie \(\ctB, \ctC\) e \(\ctD\), con \(\ctB\) e \(\ctD\) complete per idempotenti,
e un funtore \(M \colon \ctC \fun \ctD\) pieno, fedele e suriettivo a meno di retratti, il funtore indotto \(\ctB^M \colon \ctB^\ctD \fun \ctB^\ctC\)
è un'equivalenza. Con il seguente esercizio vogliamo mostrare che, nel caso in cui \(\ctB\) sia \(\ctSet\), l'ipotesi che \(\ctD\) sia
completa per idempotenti si può omettere. Per pervenire a tale scopo, piuttosto che analizzare il funtore \(\ctSet^M\) analizziamo
il suo aggiunto sinistro \(M^* \colon \ctSet^\ctC \fun \ctSet^\ctD\) introdotto nella proposizione \ref{prop_funt_indotto_pref}.
(L'esercizio è anche una scusa per mostrare un esempio di funtore iniziale.)

\begin{exercise}\label{esercizio_equiv_sutr}
 Sia \(M \colon \ctC \fun \ctD\) un funtore fra categorie piccole.
 \begin{enumerate}
  \item Per ogni oggetto \(B \in \ctSet^{\ctD}\), il funtore \(M\) induce un funtore \(M_B \colon \mathrm{El}(B \cmp M) \fun \mathrm{El}(B)\)
        che rende il seguente diagramma commutativo
        \[ \xymatrix{ \ctC^{\op} \ar[d]_{M^{\op}} & \mathrm{El}(B \cmp M) \ar[l]_-{\Phi_{B \cmp M}} \ar[d]^{M_B} \\
         \ctD^{\op} & \mathrm{El}(B) \ar[l]^{\Phi_B} }\]
  \item Se \(M\) è pieno e suriettivo a meno di retratti, allora il funtore \(M_B\) è iniziale.
  \item Se \(M\) è pieno e fedele, allora il funtore \(M^* \colon \ctSet^\ctC \fun \ctSet^\ctD\)  è pieno e fedele.
  \item Se \(M\) è pieno, fedele e suriettivo a meno di retratti, allora il funtore \(M^* \colon \ctSet^\ctC \fun \ctSet^\ctD\)
        è un'equivalenza.
 \end{enumerate}
\end{exercise}

\begin{proof}
 1. Data una freccia di \(\mathrm{El}(B \cmp M)\)
 \[ f \colon (X \in \ctC, x \in BMX) \to (Y \in \ctC, y \in BMY) \]
 con \(f \colon Y \to X\) tale che \((BMf)(y)=x\), la sua immagine attraverso \(M_B\) è la freccia
 \[ Mf \colon (MX \in \ctD, x \in BMX) \to (MY \in \ctD, y \in BMY) \]
 di \(\mathrm{El}(B)\). Con tale definizione, la commutatività del diagramma dell'enunciato è ovvia. \\
 2. Fissiamo un oggetto \((D \in \ctD, d \in BD)\) in \(\mathrm{El}(B)\).
 Per ipotesi, esiste un diagramma commutativo
 \[ \xymatrix{ D \ar[rr]^-{\id} \ar[rd]_{s} & & D \\ & MZ \ar[ru]_{r} }\]
 Otteniamo un oggetto \(r \colon (D,d) \to (MZ, (Bs)(d))\) della categoria comma \((D,d) / M_B\),
 infatti \((Br)((Bs)(d)) = (B(r \cmp s))(d) = (B\id)(d) = d\). Consideriamo ora due oggetti
 \[ \xymatrix{ & (D,d) \ar[ld]_{g} \ar[rd]^{h} \\
   (MX,x) & & (MY,y) } \]
 della categoria comma, cioè \(g \colon MX \to D\) e \(h \colon MY \to D\) tali che \((Bg)(x) = d = (Bh)(y)\).
 Possiamo costruire un nuovo diagramma commutativo
 \[ \xymatrix{ & & D \\
  MX \ar[rru]^{g} \ar[r]_-{g} & D \ar[r]_-{s} & MZ \ar[u]_{r} & D \ar[l]^-{s} & MY \ar[l]^-{h} \ar[llu]_{h} } \]
 Poiché \(M\) è pieno, esistono \(g' \colon  X \to Z\) e \(h' \colon Y \to Z\) tali che \(Mg' = s \cmp g\) e \(Mh' = s \cmp h\).
 Questo ci fornisce uno zig-zag in \((D,d) / M_B\)
 \[ \xymatrix{ & (D,d) \ar[ld]_{g} \ar[rd]^{h} \ar[d]^{r} \\
  (MX,x) & (MZ, (Bs)(d)) \ar[l]^-{Mg'} \ar[r]_-{Mh'} & (MY,y) } \]
 infatti \((BMg')(x) = (B(s \cmp g))(x) = (Bs)((Bg)(x)) = (Bs)(d)\) e, analogamente, \((BMh')(y) = (Bs)(d)\). Abbiamo così
 mostrato che \((D,d) / M_B\) è connessa, ovvero che \(M_B\) è iniziale. \\
 3.  Consideriamo il diagramma
 \[ \xymatrix{ \ctC^{\op} \ar[rr]^-{\ctY_{\ctC}} \ar[d]_{M^{\op}} \ar@{.>}[rrd]^{F} & & \ctSet^{\ctC} \ar@<-0.8ex>[d]_{M^*} \\
  \ctD^{\op} \ar[rr]_-{\ctY_{\ctD}} & & \ctSet^{\ctD} \ar@<-0.5ex>[u]_{\ctSet^M} } \]
 dove abbiamo posto \(F = \ctY_{\ctD} \cmp M^{\op}\) (vedi prova della proposizione \ref{prop_funt_indotto_pref}).
 Poiché \(F\) è pieno e fedele (perché \(M\) e \(\ctY_{\ctD}\) lo sono) e poiché, per ogni \(X \in \ctC\), il prefascio \(FX = \ctD(MX,\blank)\)
 è assolutamente presentabile, il lemma \ref{lemma_caratt_pref} garantisce che \(M^*\) è pieno e fedele. \\
 4. Assumiamo ora che ogni oggetto di \(\ctD\) sia un retratto di un  oggetto del tipo \(MZ\) con \(Z \in \ctC\). Per mostrare che
 \(M^*\) è anche essenzialmente suriettivo, consideriamo il diagramma commutativo
 \[ \xymatrix{ \mathrm{El}(B \cmp M) \ar[r]^-{\Phi_{B \cmp M}} \ar[rd]_{M_B} & \ctC^{\op} \ar[r]^-{M^{\op}}
  & \ctD^{\op} \ar[r]^-{\ctY_{\ctD}} & \ctSet^M \\
  & \mathrm{El}(B) \ar[ru]_{\Phi_B} } \]
 dove \(B \in \ctSet^{\ctD}\) e \(M_B\) è il funtore introdotto nel primo punto dell'esercizio.
 Poiché \(M_B\) è iniziale, possiamo usare il lemma \ref{lemma_caract_funt_iniz} e abbiamo
 \[ B = \colim_{\mathrm{El}(B)}(\ctY_{\ctD} \cmp \Phi_B) \simeq \colim_{\mathrm{El}(B \cmp M)}(\ctY_{\ctD} \cmp \Phi_B \cmp M_B)
  = \colim_{\mathrm{El}(B \cmp M)}(\ctY_{\ctD} \cmp M^{\op} \cmp \Phi_{B \cmp M}) \simeq \]
 \[ \simeq \colim_{\mathrm{El}(B \cmp M)}(M^* \cmp \ctY_{\ctC} \cmp \Phi_{B \cmp M}) \simeq
  M^*(\colim_{\mathrm{El}(B \cmp M)}(\ctY_{\ctC} \cmp \Phi_{B \cmp M})) = M^*(B \cmp M) \]
 Questo conclude la prova che \(M^*\) è un'equivalenza (e quindi anche \(\ctSet^M\) lo è).
\end{proof}

\section{Sottocategorie equazionali di categorie di prefasci}\label{sec_birk_prefasci}

Anticipando un problema che sarà sviluppato più ampiamente nel capitolo sulle categorie algebriche, possiamo chiederci come
riconoscere se una sottocategoria piena di \(\ctSet^\ctC\) è determinata da un insieme di equazioni su \(\ctC\) o, equivalentemente,
da una congruenza su \(\ctC\). Prima di specializzarci al caso di \(\ctSet^\ctC\), consideriamo il caso generale di \(\ctB^\ctC\).

\begin{definition}\label{def_equaz_cat}
 Sia \(\ctC\) una categoria.
 \begin{enumerate}
  \item Un'equazione su \(\ctC\), o \(\ctC\)-equazione, è una coppia \(u,v \colon X \rightrightarrows Y\) di frecce parallele di \(\ctC\).
  \item Un funtore \(M \colon \ctC \fun \ctB\) soddisfa una \(\ctC\)-equazione \((u,v)\) se \(M(u)=M(v)\).
  \item Se \(E\) è un insieme di \(\ctC\)-equazioni, denotiamo con \(\ctB^{(\ctC,E)}\) la sottocategoria piena di \(\ctB^\ctC\)
        dei funtori \(A \colon \ctC \fun \ctB\) tali che \(A(u)=A(v)\) per ogni \(\ctC\)-equazione \((u,v) \in E\). Chiamiamo le sottocategorie
        del tipo \(\ctB^{(\ctC,E)}\) sottocategorie equazionali.
 \end{enumerate}
\end{definition}

Per poter costruire la categoria quoziente, conviene rimpiazzare la nozione di \(\ctC\)-equazione con la nozione più
maneggevole di \(\ctC\)-congruenza.

\begin{definition}\label{def_cong_cat}
 Sia \(\ctC\) una categoria.
 \begin{enumerate}
  \item Una congruenza \(\sim\) su \(\ctC\), o \(\ctC\)-congruenza, è una famiglia di relazioni di equivalenza
        \[\{\sim_{X,Y} \, \subseteq \ctC(X,Y) \times \ctC(X,Y) \mid X,Y \in \ctC\}\]
        che collettivamente formano una sottocategoria della categoria prodotto \(\ctC \times \ctC\), ovvero
        tali che, se \(a \sim_{X,Y}b\) e \(c \sim_{Y,Z}d\), allora \(c \cmp a \sim_{X,Z} d \cmp b\).  \\
        (Quando questo non crea confusione, scriveremo \(a \sim b\) invece di \(a \sim_{X,Y} b\).)
  \item Se \(\sim\) è una \(\ctC\)-congruenza, denotiamo con \(\ctB^{(\ctC,\sim)}\) la sottocategoria piena di \(\ctB^{\ctC}\)
        dei funtori \(A \colon \ctC \fun \ctB\) tali che \(A(u)=A(v)\) ogniqualvolta \(u \sim v\).
 \end{enumerate}
\end{definition}

\begin{example}\label{ex_cong_funt}
 Se \(M \colon \ctC \fun \ctB\) è un funtore, otteniamo una \(\ctC\)-congruenza \(\sim_M\) nel modo seguente:
 date due frecce parallele
 \(u,v \colon X \rightrightarrows Y\) di \(\ctC\), poniamo
 \[u \sim_M v \;\mbox{ se }\; M(u)=M(v)\]
 Usando tale notazione, si ha che un funtore \(M \colon \ctC \fun \ctB\) è in \(\ctB^{(\ctC,\sim)}\) se e solo se \(\sim \,\leq\, \sim_M\).
 (La relazione d'ordine fra \(\ctC\)-congruenze è quella ovvia: \(\sim \, \leq \, \sim'\) se \(\sim_{X,Y} \, \subseteq \, \sim'_{X,Y}\)
 per ogni \(X,Y \in \ctC\).)
\end{example}

\begin{remark}\label{oss_cong_funt}
 Per completare l'esempio \ref{ex_cong_funt}, osserviamo due semplici fatti che riguardano la congruenza indotta da un funtore.
 \begin{enumerate}
  \item Dati due funtori \(M \colon \ctC \fun \ctB\) e \(N \colon \ctB \fun \ctA\), si ha
        sempre \(\sim_M \,\leq\, \sim_{N \cmp M}\).
  \item Sia \(\alpha \colon M \nat N \colon \ctC \rightrightarrows \ctB\) una trasformazione naturale. Se \(\alpha\) è un mono, allora
        \(\sim_N \,\leq\, \sim_M\). Se \(\alpha\) è un epi, allora \(\sim_M \,\leq\, \sim_N\).
 \end{enumerate}
\end{remark}

\begin{proof}
 Il punto 1 è ovvio. Per il punto 2, ricordiamoci che \(\alpha\) è un mono nella categoria \(\ctB^\ctC\) precisamente quando
 \(\alpha_X\) è un mono in \(\ctB\) per ogni oggetto \(X \in \ctC\), e analogamente per la condizione di epi (vedi corollario
 \ref{cor_fun_cat_compl}). La naturalità di \(\alpha\) applicata a due frecce \(u,v \colon X \rightrightarrows Y\) di \(\ctC\)
 dà luogo al diagramma
 \[\xymatrix{MX \ar@<-0.5ex>[rr]_-{M(v)} \ar@<0.5ex>[rr]^-{M(u)} \ar[d]_{\alpha_X} & & MY \ar[d]^{\alpha_Y} \\
  NX \ar@<-0.5ex>[rr]_-{N(v)} \ar@<0.5ex>[rr]^-{N(u)} & & NY }\]
 la cui commutatività permette di verificare le inclusioni fra \(\ctC\)-congruenze del punto 2.
\end{proof}


\begin{remark}\label{oss_inter_cong}
 Chiaramente, ogni \(\ctC\)-congruenza \(\sim\) determina un insieme \(E\) di \(\ctC\)-equazioni dato da tutte le coppie
 \((u,v)\) tali che \(u \sim v\). Vogliamo ora associare a un insieme \(E\) di \(\ctC\)-equazioni una congruenza su \(\ctC\).
 Per questo, cominciamo con l'osservare che se \(\sim\) e \(\sim'\) sono due \(\ctC\)-congruenze, la loro intersezione puntuale
 \(\sim \,\cap\, \sim'\), definita da
 \[(\sim \, \cap \, \sim' )_{X,Y} = \, \sim_{X,Y} \,\cap \, \sim'_{X,Y}\]
 è ancora una congruenza su \(\ctC\). Lo stesso dicasi per l'intersezione di una famiglia arbitraria di \(\ctC\)-congruenze.
 Ne segue che, se \(E = \{E_{X,Y} \subseteq \ctC(X,Y) \times \ctC(X,Y) \mid X,Y \in \ctC\} \) è un insieme di \(\ctC\)-equazioni,
 allora l'intersezione di tutte le \(\ctC\)-congruenze \(\sim\) tali che
 \[ E_{X,Y} \subseteq \, \sim_{X,Y} \;\mbox{ per ogni coppia }\; X,Y \in \ctC\]
 è una \(\ctC\)-congruenza ed è la più piccola \(\ctC\)-congruenza che contiene tutte le equazioni di \(E\). Denoteremo tale
 \(\ctC\)-congruenza con \(\sim_E\). Il carattere minimale di \(\sim_E\) dà luogo al seguente lemma.
\end{remark}

\begin{lemma}\label{lemma_equaz_cong}
 Sia \(E\) un insieme di \(\ctC\)-equazioni.
 Un funtore \(M \colon \ctC \fun \ctB\) soddisfa tutte le equazioni di \(E\) se e solo se \(\sim_E \,\leq\, \sim_M\).
 Ne segue che \(\ctB^{(\ctC,E)} = \ctB^{(\ctC,\sim_E)}\).
\end{lemma}

Come annunciato, a partire da una \(\ctC\)-congruenza costruiamo ora la categoria quoziente \(\ctC/{\!\sim}\) per poi
stabilirne la proprietà universale.

\begin{proposition}\label{prop_cat_quoz}
 Sia \(\ctC\) una categoria e \(\sim\) una \(\ctC\)-congruenza. Esistono una categoria \(\ctC/{\!\sim}\) e un funtore
 \(Q \colon \ctC \fun \ctC/{\!\sim}\) pieno, suriettivo sugli oggetti e tale che \(\sim \,=\,\sim_Q\). Inoltre, se
 \(M \colon \ctC \fun \ctD\) è tal che \(\sim \,\leq\, \sim_M\), allora esiste un unico \(M' \colon \ctC/{\!\sim} \fun \ctD\)
 tale che \(M' \cmp Q = M\)
 \[\xymatrix{\ctC \ar[r]^-{Q} \ar[d]_-{M} & \ctC/{\!\sim} \ar[ld]^{M'} \\ \ctD }\]
\end{proposition}

\begin{proof}
 Gli oggetti di \(\ctC/{\!\sim}\) sono gli oggetti di \(\ctC\). Le frecce di  \(\ctC/{\!\sim}\) sono le classi di equivalenza di
 frecce di \(\ctC\), cioè  \(\ctC/{\!\sim}(X,Y) = \ctC(X,Y)/\sim_{X,Y}\). Il funtore \(Q\) è definito da
 \(Q(f \colon X \to Y) = (\left[f\right] \colon X \to Y)\). Il fatto che \(Q\) sia pieno e suriettivo sugli oggetti è ovvio, così
 come la condizione \(\sim_Q \,=\, \sim\). Sia ora un funtore \(M \colon \ctC \fun \ctD\)
 tale che \(\sim \,\leq\, \sim_M\). La condizione \(M' \cmp Q = M\) impone che \(M'\) sia definito da
 \(M'(\left[f\right] \colon X \to Y) = M(f \colon X \to Y)\) e la condizione \(\sim \,\leq\, \sim_M\) implica che tale definizione
 è ben data.
\end{proof}

\begin{exercise}\label{ex_cat_quoz_bidim}
 Lasciamo a chi legge il compito di enunciare e dimostrare una proprietà universale della categoria quoziente dove la
 condizione \(M' \cmp Q = M\) sia rimpiazzata dall'esistenza di un isomorfismo naturale \(M' \cmp Q \simeq M\) e
 l'unicità di \(M'\) sia a meno di un (unico) isomorfismo naturale.
\end{exercise}

Ora possiamo dimostriare che le sottocategorie equazionali di \(\ctB^\ctC\) sono, a meno di equivalenze, le categorie
del tipo \(\ctB^{\ctC/{\!\sim}}\) per \(\sim\) una \(\ctC\)-congruenza.

\begin{corollary}\label{cor_cat_equaz_quoz}
 Sia \(\sim\) una \(\ctC\)-congruenza. Il funtore \(\ctB^Q \colon \ctB^{\ctC/{\!\sim}} \fun \ctB^\ctC\) indotto per composizione dal funtore
 universale \(Q \colon \ctC \fun \ctC/{\!\sim}\) si fattorizza attraverso la sottocategoria \(\ctB^{(\ctC,\sim)}\) e la fattorizzazione è
 un'equivalenza di categorie
 \[\xymatrix{\ctB^{\ctC/{\!\sim}} \ar[rr]^-{\ctB^Q} \ar[rd]_{\simeq} & & \ctB^{\ctC} \\ & \ctB^{(\ctC,\sim)} \ar[ru] }\]
 Ne segue che \(\ctB^{(\ctC,\sim)}\) è chiusa in \(\ctB^\ctC\) per limiti, colimiti, sottooggetti e immagini epimorfiche.
\end{corollary}

\begin{proof}
 Se \(B \in \ctB^{\ctC/{\!\sim}}\), allora \(\ctB^Q(B) = B \cmp Q \in \ctB^{(\ctC,\sim)}\) perché \(\sim \,=\, \sim_Q \;\leq\, \sim_{B \cmp Q}\).
 Verifichiamo ora che la corestrizione di \(\ctB^Q\) a \(\ctB^{(\ctC,\sim)}\) è un'equivalenza. Se
 \(\beta \colon B_1 \nat B_2 \colon \ctC/{\!\sim} \rightrightarrows \ctB\) è una trasformazione naturale, allora
 \((\ctB^Q(\beta))_X = \beta_{QX} = \beta_X\) per ogni \(X \in \ctC\). La fedeltà di \(\ctB^Q\) è dunque ovvia. Per quanto riguarda
 la pienezza, se \(\alpha \colon B_1 \cmp Q \nat B_2 \cmp Q \colon \ctC \rightrightarrows \ctB\) è una trasformazione naturale,
 otteniamo una trasformazione naturale \(\beta \colon B_1 \nat B_2 \colon \ctC/{\!\sim} \rightrightarrows \ctB\)
 ponendo \(\beta_X = \alpha_X \colon B_1X = B_1(QX) \to B_2(QX) = B_2X\) per ogni \(X \in \ctC\). Per la naturalità di \(\beta\),
 basta osservare che se \(\left[u\right] \colon X \to Y\) è una freccia di \(\ctC/{\!\sim}\), allora
 \[B_2\left[u\right] \cmp \beta_X = B_2(Qu) \cmp \alpha_X = \alpha_Y \cmp B_1(Qu) = \beta_Y \cmp B_1\left[u\right] \]
 Per il carattere essenzialmente suriettivo, se \(A \in \ctB^{(\ctC,\sim)}\), allora \(\sim \,\leq\, \sim_A\) e quindi la proprietà
 universale \ref{prop_cat_quoz} garantisce l'esistenza di un (unico) \(A' \in \ctB^{\ctC/{\!\sim}}\) tale che \(\ctB^Q(A')=A\). \\
 Le condizioni di chiusura di \(\ctB^{(\ctC,\sim)}\) seguono ora dal corollario \ref{cor_funt_indotto_catfunt} (per i limiti e i colimiti) e
 dall'osservazione \ref{oss_cong_funt} (per i sottooggetti e le immagini epimorfiche).
\end{proof}

Nel prossimo lemma caratterizziamo le categorie quoziente. Si tratta di una caratterizzazione a meno di equivalenza di
categorie, ragione per cui la condizione su \(Q\) che metteremo in evidenza è di essere essenzialmente suriettivo sugli
oggetti, non strettamente suriettivo.

\begin{lemma}\label{lemma_caratt_categ_quoz}
 Sia \(M \colon \ctC \fun \ctD\) un funtore e sia \(\sim_M\) la \(\ctC\)-conguenza indotta de \(M\). Consideriamo la situazione
 della proprietà universale \ref{prop_cat_quoz}:
 \[\xymatrix{\ctC \ar[r]^-{Q} \ar[d]_-{M} & \ctC/{\!\sim_M} \ar[ld]^{M'} \\ \ctD }\]
 Il funtore \(M'\) è sempre fedele. Inoltre \(M'\) è pieno se e solo se \(M\) lo è ed è essenzialmente suriettivo se e solo
 se \(M\) lo è. Ne segue che \(M'\) è un'equivalenza se e solo se \(M\) è pieno ed essenzialemente suriettivo.
\end{lemma}

\begin{proof}
 Verifichiamo per esempio che \(M'\) è fedele:
 \[M'\left[f\right] = M'\left[g\right] \;\Rightarrow\; M(f) = M(g) \;\Rightarrow\; f \sim_M g \;\Rightarrow\; \left[f\right] = \left[g\right]\]
 Invitiamo chi legge a terminare la dimostrazione come esercizio.
\end{proof}

%\begin{corollary}\label{cor_ind_plfed}
%Se un funtore \(M \colon \ctC \fun \ctD\) è pieno ed essenzialmente suriettivo, il funtore \(\ctB^M \colon \ctB^D \fun \ctB^\ctC\) 
%indotto per composizione è pieno e fedele.
%\end{corollary} 
%
%\begin{proof}
%Con la notazione del lemma \ref{lemma_caratt_categ_quoz}, \(M = M' \cmp Q\). Ne segue che \(\ctB^M = \ctB^Q \cmp \ctB^{M'}\) 
%con \(\ctB^{M'}\) un'equivalenza per il lemma \ref{lemma_caratt_categ_quoz} e \(\ctB^Q\) pieno e fedele per il corollario 
%\ref{cor_cat_equaz_quoz}.
%\end{proof} 

\begin{remark}\label{oss_restr_quot_pref}
 Per migliorare il corollario \ref{cor_cat_equaz_quoz}, ci restringiamo adesso al caso in cui la categoria \(\ctB\) che
 compare nel lemma \ref{lemma_equaz_cong} e nel corollario \ref{cor_cat_equaz_quoz} sia la
 categoria \(\ctSet\). Per un insieme \(E\) di \(\ctC\)-equazioni abbiamo quindi un'equivalenza
 \(\ctSet^{(\ctC,E)} = \ctSet^{(\ctC,\sim_E)} \simeq \ctSet^{\ctC/{\!\sim_E}}\). Inoltre, a meno di tale equivalenza, l'inclusione
 \(\ctSet^{(\ctC,E)} \fun \ctSet^\ctC\) coincide con il funtore \(\ctSet^Q \colon \ctSet^{\ctC/{\!\sim_E}} \fun \ctSet^\ctC\) che
 sappiamo preservare i colimiti e avere un aggiunto sinistro \(Q^*\), vedi proposizione \ref{prop_funt_indotto_pref}.
 Ricordiamo anche che per le categorie del tipo \(\ctSet^\ctC\) non c'è differenza fra epi ed epi regolare (esempio
 \ref{esempio_set_reg}).
\end{remark}

\begin{lemma}\label{lemma_unita_agg_ind}
 Sia \(M \colon \ctC \fun \ctD\) un funtore fra categorie piccole e sia
 \[\xymatrix{ \ctSet^\ctD \ar@<-0.5ex>[rr]_-{\ctSet^M} & & \ctSet^\ctC \ar@<-0.5ex>[ll]_-{M^*} } \;\;\;\;\;\; M^* \dashv \ctSet^M\]
 l'aggiunzione indotta. Se \(M\) è pieno, allopra l'unità \(\eta^A \colon A \to \ctSet^M(M^*A)\) dell'aggiunzione è un
 epimorfismo (regolare) per ogni \(A \in \ctSet^\ctC\).
\end{lemma}

\begin{proof}
 Ricordiamo la situazione della proposizione \ref{prop_funt_indotto_pref}
 \[ \xymatrix{\ctC^{\op} \ar[d]_{M^{\op}} \ar[r]^-{\ctY_{\ctC}} & \ctSet^{\ctC} \ar@<-0.5ex>[d]_{M^*} \\
  \ctD^{\op} \ar[r]_-{\ctY_{\ctD}} & \ctSet^{\ctD} \ar@<-0.5ex>[u]_{\ctSet^M} } \]
 con \(M^* \dashv \ctSet^M\) e \(\ctSet^M\) che preserva i colimiti. Se \(A \in \ctSet^\ctC\), sappiamo che
 \[A \simeq \colim_{\Elts\ctD A}(\xymatrix{\Elts\ctD A \ar[r]^-{\Phi_A} & \ctC^{\op} \ar[r]^-{\ctY_\ctC} & \ctSet^\ctC })\]
 (vedi proposizione \ref{prop_colim_rappr}), da cui segue che
 \[\begin{array}{rcl} \ctSet^M(M^*A)                                              & \simeq                         &
             \colim_{\Elts\ctD A}(\xymatrix{\Elts\ctD A \ar[r]^-{\Phi_A} & \ctC^{\op} \ar[r]^-{\ctY_\ctC} & \ctSet^\ctC \ar[r]^-{M^*}      &
             \ctSet^\ctD \ar[r]^-{\ctSet^M}                                    & \ctSet^\ctC})                                                     \\
                                                                               & \simeq                         &
             \colim_{\Elts\ctD A}(\xymatrix{\Elts\ctD A \ar[r]^-{\Phi_A} & \ctC^{\op} \ar[r]^-{M^{\op}}   & \ctD^{\op} \ar[r]^-{\ctY_\ctD} &
             \ctSet^\ctD \ar[r]^-{\ctSet^M}                                    & \ctSet^\ctC})
  \end{array}\]
 Possiamo dunque descrivere l'unità dell'aggiunzione tramite il diagramma seguente
 \[\xymatrix{A \ar[rr]^-{\eta^A} & &  \ctSet^M(M^*A) \\
  \colim_{\Elts\ctD A} \ctC(X,\blank) \ar[u]^{\simeq} & & \colim_{\Elts\ctD A}(\ctD(MX,\blank) \cmp M) \ar[u]_{\simeq} \\
  \ctC(X,\blank) \ar[u]^{\sigma^{(X,x)}} \ar[rr]_-{\eta^{\ctC(X,\blank)}} & & \ctD(MX,\blank) \cmp M \ar[u]_{\sigma^{(X,x)}} }\]
 Poiché in \(\ctSet^\ctC\) i colimiti si calcolano punto a punto, il diagramma precedente induce, per ogni \(Y \in \ctC\),
 un diagramma commutativo
 \[\xymatrix{\colim_{\Elts\ctD A}\ctC(X,Y) \ar[rr]^-{\eta^A_Y} & & \colim_{\Elts\ctD A}\ctD(MX,MY) \\
  \ctC(X,Y) \ar[u]^{\sigma^{(X,x)}_Y} \ar[rr]_-{\eta^{\ctC(X,\blank)}_Y} & & \ctD(MX,MY) \ar[u]_{\sigma^{(X,x)}_Y} }\]
 Rimane da descrivere la componente in \(Y\) della trasformazione naturale \(\eta^{\ctC(X,\blank)}\), ma questa è
 data semplicemente dall'azione di \(M\) sulle frecce di \(\ctC\):
 \[\eta^{\ctC(X,\blank)}_Y = M_{X,Y} \colon \ctC(X,Y) \to \ctD(MX,MY)\]
 Poiché, per ipotesi, \(M\) è pieno, tale funzione \(M_{X,Y}\) è suriettiva. Ne segue che anche la funzione
 \(\eta^A_Y\) è suriettiva e, poiché quasto vale per ogni \(Y \in \ctC\), possiamo concludere che \(\eta^A\) è
 un epi.
\end{proof}

Grazie al lemma \ref{lemma_unita_agg_ind} otteniamo la seguente caratterizzazione delle sottocategorie equazionali delle
categorie di prefasci. Per la nozione di sottocategoria	 epiriflessiva, che appare nella prossima proposizione, si veda la
definizione \ref{def_sottocat_epirifl} tenendo presente quanto ricordato nell'osservazione \ref{oss_restr_quot_pref} a proposito
degli epi regolari nelle categorie di prefasci.

\begin{proposition}\label{prop_stc_equaz_pref}
 Siano \(\ctC\) una categoria piccola e \(\ctA\) una sottocategoria piena e repleta di \(\ctSet^\ctC\).
 Le condizioni seguenti sono equivalenti:
 \begin{enumerate}
  \item \(\ctA\) è equivalente a \(\ctSet^{(\ctC,E)}\) per un qualche insieme \(E\) di \(\ctC\)-equazioni.
  \item \(\ctA\) è epiriflessiva e chiusa per colimiti in \(\ctSet^\ctC\).
  \item \(\ctA\) è chiusa in \(\ctSet^\ctC\) per colimiti, prodotti e sottooggetti.
 \end{enumerate}
\end{proposition}

\begin{proof}
 \(1 \Rightarrow 2\): Poiché \(\ctSet^{(\ctC,E)} = \ctSet^{(\ctC,\sim_E)}\), possiamo considerare il funtore universale
 \(Q \colon \ctC \fun \ctC/{\!\sim_E}\) della proposizione \ref{prop_cat_quoz} e il diagramma commutativo
 \[\xymatrix{\ctSet^{\ctC/{\!\sim_E}} \ar[rr]^-{\ctSet^Q} \ar[rd]_{\simeq} & & \ctSet^{\ctC} \\ & \ctSet^{(\ctC,\sim_E)} \ar[ru] }\]
 del corollario \ref{cor_cat_equaz_quoz}. Ne deduciamo che \(\ctSet^{(\ctC,\sim_E)}\) è chiusa in \(\ctSet^\ctC\)
 per colimiti (perché \(\ctSet^Q\) preserva i colimiti), che è riflessiva (perché \(\ctSet^Q\) ha aggiunto sinistro) e,
 in particolare, epiriflessiva per il lemma \ref{lemma_unita_agg_ind} (perché \(Q\) è pieno). \\
 \(2 \Rightarrow 1\): Sia \(R \colon \ctSet^\ctC \fun \ctA\) l'aggiunto sinistro all'inclusione \(i \colon \ctA \fun \ctSet^\ctC\).
 Come nella dimostrazione del corollario \ref{cor_stab_pref}, poniamo
 \[\ctG = \{R(\ctC(X,\blank)) \mid X \in \ctC\}\]
 che guardiamo come sottocategoria piena di \(\ctA\). Sappiamo già che \(\ctA \simeq \ctSet^{\ctG^{\op}}\) (vedi proposizione
 \ref{prop_caratt_pref}). Definiamo ora un funtore \(M \colon \ctC \fun \ctG^{\op}\) chiedendo che \(M^{\op} \colon \ctC^{\op} \fun \ctG\)
 sia la corestizione di \(R \cmp \ctY_\ctC\) a \(\ctG\)
 \[\xymatrix{\ctC^{\op} \ar[rr]^-{\ctY_\ctC} \ar[d]_{M^{\op}} & & \ctSet^\ctC \ar[d]^{R} \\
  \ctG \ar[r]_-{\ctY_{\ctG^{\op}}} & \ctSet^{\ctG^{\op}} \ar[r]_-{\simeq} & \ctA }\]
 Tale funtore \(M\) è ovviamente suriettivo sugli oggetti (per definizione di \(\ctG\)).
 Se mostriamo che è anche pieno, allora il lemma \ref{lemma_caratt_categ_quoz} garantisce che \(\ctG^{\op} \simeq \ctC/{\!\sim_M}\)
 e quindi \(\ctA \simeq \ctSet^{\ctG^{\op}} \simeq \ctSet^{\ctC/{\!\sim_M}}\), il che conclude la prova di questa implicazione. Sia dunque
 \[g \colon MX = R(\ctC(X,\blank)) \to R(\ctC(Y,\blank))= MY\]
 una freccia in \(\ctG^{\op}\). Nell'aggiunzione \(R \dashv i\), tale freccia corrisponde a una freccia
 \[ g' \colon \ctC(X,\blank) \to i(R(\ctC(Y,\blank)))\]
 Poiché l'unità \(\eta^{\ctC(Y,\blank)} \colon \ctC(Y,\blank) \to i(R(\ctC(Y,\blank)))\) è un epi e \(\ctC(X,\blank)\) è un oggetto proiettivo, esiste una freccia
 \(g'' \colon \ctC(X,\blank) \to \ctC(Y,\blank)\) tale che \(\eta^{\ctC(Y,\blank)} \cmp g'' = g'\). Per il lemma di Yoneda, la freccia \(g''\) corrisponde a sua volta
 a una feeccia \(f \colon Y \to X\) in \(\ctC\) e si verifica senza problemi che \(M(f)=g\). \\
 \(1 \Rightarrow 3\): È un caso particolare del corollario \ref{cor_cat_equaz_quoz}. \\
 \(3 \Rightarrow 2\): Poiché \(\ctSet^\ctC\) è completa, regolare (vedi esempio \ref{esempio_set_reg}) e co-well-powered (vedi esempio
 \ref{esempio_set_cowellpow}), possiamo applicare il lemma \ref{lemma_caract_epirifl} e concludere che \(\ctA\) è epiriflessiva.
\end{proof}

\color{black}